\chapter{Sequence design and graph encoding}
Suppose four vertices have four different DNA words. Does that make every
directed edge chemically distinguishable? No. Two words may differ at every
aligned position and still offer a long complementary segment after one strand
slides. Two components can create a new recognition word at their junction.
Symbolic uniqueness and molecular selectivity are different contracts.

We will build a graph codec, prove what it preserves, and challenge its sequence
screens. The eight-base words are teaching objects, not synthesis recommendations.
Chapters 5--7 supply polarity, equilibrium and kinetics; Chapter 10 explains
why a measured band cannot by itself establish molecular identity.

\section{Choose a representation before measuring distance}
\index{reverse complement}\index{strand polarity}
Every sequence string is written $5'$ to $3'$. Let $c$ exchange A with T
and C with G. For $s=s_0s_1\cdots s_{L-1}$, define
\begin{equation}
 \operatorname{rc}(s)_i=c(s_{L-1-i}),\qquad 0\le i<L.
\end{equation}
Complementation chooses the paired base; reversal puts the opposing strand
back into the same written orientation. Since $c(c(b))=b$, applying the
reverse complement twice returns the original word. Length is unchanged.
\Plate{01-polarity}{The physical duplex has antiparallel backbones, while
both stored strings use the $5'$-to-$3'$ convention. Reversing the lower
strand's written description reconciles the two views. Dotted rungs indicate
pairing, not backbone bonds.}{fig:polarity}
\Listing{revcomp}
The companion helper \texttt{dna} rejects empty strings, lowercase letters
and ambiguous bases. This is a reference-code contract, not a claim that
IUPAC ambiguity codes are invalid biology. Degenerate bases require different
equality and pairing rules. Tests check the involution on all 256 four-base words.

For equal-length strings, the Hamming distance is
\begin{equation}
 d_H(a,b)=\sum_{i=0}^{L-1}\mathbf1[a_i\ne b_i].
\end{equation}
It counts substitutions in a fixed alignment. It is neither a binding energy
nor an edit distance.
\Listing{hamming}

\section{A distance guarantee has a noise model}
\index{Hamming distance}\index{error correction}
Let a codebook have minimum pairwise distance $d_{\min}$. Suppose a read has
at most $t$ substitutions, with unchanged length and known reading frame.
Nearest-codeword decoding is guaranteed to recover the original if
$d_{\min}>2t$. To prove this, let $a$ be the original, $y$ the observation
and $b$ another codeword. The triangle inequality gives
\begin{equation}
 d_H(y,b)\ge d_H(a,b)-d_H(a,y)\ge d_{\min}-t>t\ge d_H(y,a).
\end{equation}
At $d_{\min}=2t$, two words can tie; an arbitrary tie-break does not restore
evidence. The proof excludes indels, mixed reads and hybridization. Our exact
decoder below does not implement error correction: a theorem in the text
does not add a feature to the program.

Error tolerance consumes sequence space. A radius-$t$ Hamming ball around
one length-$L$ word contains
\begin{equation}
 V_4(L,t)=\sum_{i=0}^t \binom{L}{i}3^i
\end{equation}
words: choose the changed positions, then three alternatives at each.
Disjoint balls for $M$ codewords require $M V_4(L,t)\le4^L$.
At $L=8,t=1$, $V_4=25$, giving $M\le2621$. This is an upper bound,
not a construction. Folding and composition constraints can reduce the
usable set further. Longer words buy representation space while changing
the molecules whose chemistry must be evaluated.

\section{Slide the alignment}
\index{off-target interaction}
Take $a=\texttt{ACGTAC}$ and $b=\texttt{TACGTA}$. Their Hamming distance
is six, the maximum. But $b=\operatorname{rc}(b)$ and
$a[0{:}5]=\operatorname{rc}(b)[1{:}6]=\texttt{ACGTA}$.
Half-open slices include the starting index and exclude the ending index.
There is a five-base complementary opportunity at an offset.
\Plate{02-offset}{The fixed same-direction comparison has six mismatches.
Shifting the actual opposing strand exposes five complementary positions.
Unpaired ends protrude. This geometry is a warning, not a binding-yield
prediction.}{fig:offset}
Let $D_{ij}$ be the length of the common suffix ending at $a_i,b_j$:
\begin{equation}
 D_{ij}=\begin{cases}1+D_{i-1,j-1},&a_i=b_j,\\0,&a_i\ne b_j.\end{cases}
\end{equation}
Boundary values are zero. The largest cell is the longest common contiguous
word. Comparing $a$ with $\operatorname{rc}(b)$ makes this a complementary
segment screen.
\Listing{longest_match}
Only the previous row is needed: time $O(|a||b|)$, storage $O(|b|)$.
The returned start positions expose the actual alignment. Ties retain
the first maximum in loop order. Tests compare the dynamic program with
explicit substring enumeration on a small alphabet.

An exact-run score omits interrupted stems, bulges, loop penalties,
concentrations and competing complexes. A low score does not certify zero
crosstalk, and a short exact segment need not be stable under the intended
conditions. Use the score to reject obvious conflicts or prioritize analysis,
not to replace thermodynamics with an integer.

\section{A strand can compete with its own partner}
\index{hairpin}\index{secondary structure}
In \texttt{GCGAAACGC}, the first three bases are reverse-complementary to the
last three. They admit a stem separated by a three-base loop. A pairing
drawing specifies a possible structure, not its probability.
\Plate{03-hairpin}{An exact inverted repeat admits this hairpin-like pattern.
One continuous backbone runs around the loop; the dotted stem rungs are
base pairs. Atomic structure and solution populations are not represented.}{fig:hairpin}
For stem length $k$ and minimum loop length $\ell$, arm starts $i,j$ must
satisfy $j\ge i+k+\ell$. The screen tests
\begin{equation}
 s[i{:}i+k]=\operatorname{rc}(s[j{:}j+k]).
\end{equation}
The separation rule prevents one base from occupying both arms. Increasing
the minimum loop length from three to four removes this particular warning,
not every possible secondary structure.
\Listing{hairpins}
NUPACK distinguishes sequence constraints, structural ensembles and tube-level
design, including intended complexes and competing species~\cite{nupack}.
Our code does not invoke or reproduce that package. A subsequent design stage
needs complete structures, concentrations, physical conditions and an appropriate
thermodynamic model, followed by experimental validation.

\section{Construct a small inspectable library}
\index{sequence design}
Choose three equal-length words from a supplied list. Require GC fraction
between $0.375$ and $0.625$, no three-base homopolymer, no exact three-base
hairpin warning with minimum loop three, and self reverse-complement Hamming
distance at least three. Each distinct pair must have direct and
reverse-complement distances at least three and a longest complementary
run at most four.

These are declared screens, not universal safety thresholds. Saying a library
is orthogonal is incomplete unless the excluded interactions are specified.
The reference exhaustively checks small subsets and returns the first feasible
one in lexical order. It neither minimizes free energy nor maximizes library size.
\Listing{choose_library}
\begin{center}\input{\Generated/selected-table.tex}\end{center}
The twelve-candidate example gives the displayed three words. If no subset
works, the function returns \texttt{None}; it does not silently relax a
constraint. With $N$ surviving candidates, choosing $M$ words can require
$\binom{N}{M}$ subset checks. This is a transparent small-instance reference,
not an industrial search algorithm. A faster solver must still certify the
same stated predicates.

\section{Expose a shared vertex at an edge junction}
\index{graph encoding}\index{splint}\index{ligation}
Use four separately specified vertex words:
\begin{center}\input{\Generated/code-table.tex}\end{center}
These graph-demo words are not asserted to satisfy every screen above.
Split $C(v)=L_vR_v$ into four-base halves and represent $u\to v$ by
\begin{equation}
 E(u,v)=R_uL_v.
\end{equation}
The graph contains $A\to B,A\to C,B\to C,C\to B,B\to D,C\to D$.
\begin{center}\input{\Generated/edge-table.tex}\end{center}
Consecutive edges expose $R_u\,L_v\,R_v\,L_w$. A complement of $L_vR_v$
can align their middle domains across a nick. Adleman used overlapping
edge oligonucleotides and vertex splints, with special endpoint handling
in the original experiment~\cite{adleman}. Our codec instead uses explicit
digital terminal halves; it is not that laboratory recipe.
\Plate{04-graph}{At B, consecutive edge strands expose CCTA and GTCA.
The lower B splint pairs antiparallel across the nick. Graph arrows specify
allowed traversal. A suitable ligation reaction can join the upper backbone;
hybridization alone does not make that covalent bond.}{fig:graph}

Distinct whole vertex words do not guarantee distinct edge words. If two
sources share the same right half, their edges to a fixed destination collide.
Our stronger codec contract requires every half-domain to be distinct.
\Listing{encode_edges}
Then $E(u,v)=E(u',v')$ implies $R_u=R_{u'}$ and $L_v=L_{v'}$,
hence $u=u'$ and $v=v'$. This injectivity proof concerns exact strings.
It implies neither zero partial binding nor equal reaction rates.

\section{Round-trip the walk, then attack the boundaries}
For a walk $(v_0,\ldots,v_k)$, append the leading $L_{v_0}$ and trailing
$R_{v_k}$ to the edge words:
\begin{align}
 L_{v_0}\prod_{i=0}^{k-1}E(v_i,v_{i+1})R_{v_k}
 &=L_{v_0}R_{v_0}L_{v_1}R_{v_1}\cdots L_{v_k}R_{v_k}\\
 &=C(v_0)C(v_1)\cdots C(v_k).
\end{align}
The product denotes ordered concatenation, not concentration multiplication.
There are $8(k+1)$ bases. Repeated vertices are allowed: a walk is not
automatically Hamiltonian.
\Listing{assemble}
\Listing{decode}
Encoding then decoding $A\to B\to C\to D$ returns \texttt{ABCD}.
Tests enumerate all valid walks through five vertex visits and compare
against direct vertex-word concatenation. This second oracle matters:
an encoder and decoder could share the same mistake.

Decoding assumes a known frame and exact words. One insertion can destroy
every subsequent boundary. Reverse-complement reads need an explicit
orientation policy. Neither failure is fixed by the distance proof for
substitutions.

Finally, concatenate \texttt{AACG} and \texttt{TCAG}. Neither part contains
\texttt{CGTC}, but the composed sequence does, starting at index two.
\Plate{05-seam}{A recognition word appears across a domain boundary.
Screening isolated parts misses the highlighted window. Whole-construct
validation must include junctions and terminal handles.}{fig:seam}
The resulting design discipline has two levels: prove representation semantics,
then examine the actual molecular constructs, including unintended interactions.
Chapter 12 asks how finite candidate populations and imperfect filters affect
the ability to recover a valid encoded witness.

\section{Exercises with worked reasoning}
\begin{enumerate}
\item \textbf{Orientation.} Find $\operatorname{rc}(\texttt{AACG})$.
 \emph{Solution:} complement TTGC, then reverse CGTT.
\item \textbf{Correction.} Distance for two substitutions?
 \emph{Solution:} at least five, with fixed length and known frame.
\item \textbf{Packing.} Evaluate $V_4(8,2)$.
 \emph{Solution:} $1+24+28\cdot9=277$, so $M\le236$; not a construction.
\item \textbf{Offset.} Explain the six-mismatch failure.
 \emph{Solution:} an offset reveals ACGTA in the reverse-complement comparison.
 Binding energetics remain uncomputed.
\item \textbf{Recurrence.} Why reset on mismatch?
 \emph{Solution:} the common suffix must be contiguous.
\item \textbf{Loop.} For length nine, $k=3,\ell=3$, which arm starts fit?
 \emph{Solution:} only $(0,6)$.
\item \textbf{Collision.} Give distinct sources that collide.
 \emph{Solution:} AACGTCAG and CCTATCAG share their right half.
\item \textbf{Direction.} Compute $E(B,C)$ and $E(C,B)$.
 \emph{Solution:} GTCAGATC and TGACCCTA.
\item \textbf{Walk length.} Encode a five-visit walk with repetition.
 \emph{Solution:} ABCBD is valid and has 40 bases; it is not Hamiltonian.
\item \textbf{Testing.} How can round-trip testing pass a wrong codec?
 \emph{Solution:} shared convention bugs; add literal word and forbidden-edge oracles.
\item \textbf{Composition.} Locate the seam-spanning word.
 \emph{Solution:} AACGTCAG$[2{:}6]=$CGTC.
\item \textbf{Evidence.} Design the next validation stage.
 \emph{Rubric:} specify complete constructs, intended/competing complexes,
 concentrations and conditions, thermodynamic and kinetic limits, controls,
 and an identity-discriminating readout.
\end{enumerate}
\section*{Selective glossary}
\addcontentsline{toc}{section}{Selective glossary}
\textbf{Codeword:} fixed string assigned to an object.
\textbf{Orthogonality:} exclusion of specified interactions.
\textbf{Splint:} complementary strand aligning other strands.
\textbf{Reading frame:} origin and width of sequence partitioning.
\textbf{Injective encoding:} distinct objects have distinct encoded strings.
